% ============================================================================
% Fichier  : partie5/P05_C21_detection.tex
% Titre    : Se Faire Détecter : le Point de Vue du Défenseur
% Auteur   : MINKA MI NGUIDJOÏ Thierry Emmanuel
% Version  : 1.0 -- novembre 2026
% ============================================================================
% Sources : MITRE ATT&CK ; journaux de sécurité Windows ; principes de
% détection. Ce chapitre prend le point de vue défensif : il n'enseigne
% aucune technique d'évasion. Opérations LOGONE et MBAM : cas fictifs.
% ============================================================================

\chapter{Se Faire Détecter : le Point de Vue du Défenseur}
\label{chap:detection}
\chaptermark{Le point de vue du défenseur}

\epigraph{%
    « Le meilleur testeur offensif est celui qui apprend au défenseur\\
    à le voir la prochaine fois. »%
}{Principe du cours}

\niveauChapitre{Intermédiaire à expert}{Chapitre~\ref{chap:rouge}~;
Chapitre~\ref{chap:post}.}

\begin{boxobjectifs}
\begin{itemize}
    \item Comprendre comment un défenseur détecte chaque étape d'une
          attaque.
    \item Relier les techniques offensives étudiées à leurs traces.
    \item Construire une matrice de couverture de détection.
    \item Écrire une règle de détection à partir d'une technique jouée.
    \item Faire du testeur offensif un contributeur à la défense.
\end{itemize}
\end{boxobjectifs}

\bigskip

\begin{boiteRemarque}{Orientation défensive}
Ce chapitre prend délibérément le point de vue du défenseur. Il n'enseigne
pas à contourner la détection~: il enseigne à la mesurer et à l'améliorer.
Le testeur éthique joue ses techniques \emph{au vu et au su} du défenseur,
pour renforcer ce dernier.
\end{boiteRemarque}

% ============================================================================
\section*{Question directrice : pour chaque technique jouée, le défenseur disposait-il d'une trace et d'une règle~?}
% ============================================================================

Dans l'Opération MBAM, 30~techniques ont été jouées, 18~détectées, dont la
moitié seulement en direct. La question de ce chapitre est systématique~:
pour chacune des techniques étudiées dans ce cours, quelle est sa trace, et
quelle règle la transforme en alerte~?

\medskip
\textbf{Plan du chapitre.}
\begin{enumerate}
    \item Voir chaque étape (section~\ref{sec:c21-etapes}).
    \item Des techniques aux traces (section~\ref{sec:c21-traces}).
    \item La matrice de couverture (section~\ref{sec:c21-matrice}).
    \item Écrire une règle (section~\ref{sec:c21-regle}).
    \item Le testeur, allié du défenseur (section~\ref{sec:c21-allie}).
\end{enumerate}

\separateur

% ============================================================================
\section{Voir Chaque Étape}
\label{sec:c21-etapes}
% ============================================================================

\begin{tableaucours}{L{3.6cm}L{4.4cm}L{4.2cm}}{%
    \tableauHeader{Étape} &
    \tableauHeader{Trace typique} &
    \tableauHeader{Source}
}
Balayage & Nombreuses connexions refusées & Pare-feu, flux réseau \\
Hameçonnage & Courriel signalé, clic sur lien connu & Messagerie, mandataire
web \\
Authentification suspecte & Connexion inhabituelle, NTLM depuis un poste
anormal & Journaux du domaine \\
Mouvement latéral & Service distant installé, session réseau & Événements
7045, 4624 \\
Persistance & Compte créé, tâche planifiée & Événements 4720, 4698 \\
\end{tableaucours}
\vspace{-0.8em}
\begin{center}{\small\itshape Tableau~21.1. Voir chaque étape d'une attaque}\end{center}

% ============================================================================
\section{Des Techniques aux Traces}
\label{sec:c21-traces}
% ============================================================================

Chaque technique des Parties~III à~V laisse une trace, déjà signalée dans
son chapitre. Le défenseur regroupe ces traces par source de journal et
s'assure que chacune est bien collectée~: une trace non journalisée est une
technique invisible.

% ============================================================================
\section{La Matrice de Couverture}
\label{sec:c21-matrice}
% ============================================================================

\begin{boxdef}
\termecle{Matrice de couverture de détection}~: tableau qui croise les
techniques (par exemple celles d'ATT\&CK) avec l'état de la détection~: trace
collectée, règle existante, règle testée. Elle montre les trous.
\end{boxdef}

\begin{boxprocedure}
\textbf{Construire une matrice de couverture}
\begin{enumerate}
    \item Lister les techniques jouées ou attendues.
    \item Pour chacune, noter la trace et sa source.
    \item Indiquer si la trace est collectée.
    \item Indiquer si une règle existe, et si elle a été testée.
    \item Prioriser les trous selon la gravité de la technique.
\end{enumerate}
\end{boxprocedure}

% ============================================================================
\section{Écrire une Règle}
\label{sec:c21-regle}
% ============================================================================

Une règle de détection s'exprime d'abord en français~: \emph{quelle
condition, sur quelle source, déclenche quelle alerte}. On la teste en
rejouant la technique (phase violette), on mesure les fausses alertes, puis
on l'ajuste. Le lab HE2-LAB-09 fait écrire et tester de telles règles.

\begin{boxscenario}{Opération LOGONE. Les règles qui manquaient}
\textbf{Constatation.} Les traces de la connexion anormale du compte de la
facturation et de son ajout à un groupe à privilèges existaient, mais aucune
règle ne les transformait en alerte.

\medskip
\textbf{Analyse} (confiance élevée). Après restitution, deux règles ont été
écrites et testées~: « compte à privilèges se connectant depuis un poste
ordinaire » et « ajout d'un compte à un groupe sensible hors fenêtre de
maintenance ». La détection ne coûtait pas d'outil nouveau, seulement des
règles.
\end{boxscenario}

% ============================================================================
\section{Le Testeur, Allié du Défenseur}
\label{sec:c21-allie}
% ============================================================================

\begin{boxmetier}{Analyste SOC / Red Team}
\textbf{Livrer les traces, pas seulement les failles.} À la fin d'une
mission, l'analyste remet au défenseur la liste horodatée de ses actions~:
c'est le meilleur jeu de test pour vérifier et compléter la détection.
\end{boxmetier}

\begin{boxalert}
\textbf{Violation typique.} Se vanter d'être « resté invisible » sans aider
le défenseur à le devenir moins.

\smallskip
\textbf{Règle~:} rester discret est un moyen de mesurer la détection, pas une
fin. Le testeur éthique livre toujours ce qu'il faut pour que le défenseur
le détecte la fois suivante.
\end{boxalert}

\separateur

\begin{boxresume}
\textbf{Ce qu'il faut retenir de ce chapitre~:}
\begin{itemize}
    \item \textbf{Chaque étape} d'une attaque laisse une trace, dans une
          source précise.
    \item \textbf{Une trace non collectée} rend la technique invisible.
    \item \textbf{La matrice de couverture} montre les trous de détection.
    \item \textbf{Une règle} s'énonce en français, se teste, s'ajuste.
    \item \textbf{Le testeur} livre ses actions horodatées pour nourrir la
          défense.
\end{itemize}
\end{boxresume}

\begin{boxinfo}
\textbf{Points de contrôle de l'élève.} Je sais~: relier une étape à sa
trace et à sa source~; expliquer pourquoi une trace doit être collectée~;
construire une matrice de couverture~; énoncer une règle de détection~;
décrire l'apport du testeur à la défense.
\end{boxinfo}

\bigskip

\begin{boxexo}[title={\ding{45}\quad Exercice~21.1. Traces \niveaubase}]
Donnez la trace et la source de~: un balayage~; la création d'un compte.
\end{boxexo}

\begin{boxexo}[title={\ding{45}\quad Exercice~21.2. Invisible \niveaubase}]
Pourquoi une technique dont la trace n'est pas collectée est-elle
invisible~?
\end{boxexo}

\begin{boxexo}[title={\ding{45}\quad Exercice~21.3. Matrice \niveaubase}]
Quelles colonnes comporte une matrice de couverture de détection~?
\end{boxexo}

\begin{boxexo}[title={\ding{45}\quad Exercice~21.4. Règle en français \niveaumoyen}]
Énoncez en français la règle qui détecte un compte à privilèges se
connectant depuis un poste ordinaire.
\end{boxexo}

\begin{boxexo}[title={\ding{45}\quad Exercice~21.5. Fausses alertes \niveaumoyen}]
Pourquoi faut-il mesurer les fausses alertes d'une règle avant de la
déployer~?
\end{boxexo}

\begin{boxexo}[title={\ding{45}\quad Exercice~21.6. En direct ou après coup \niveaumoyen}]
Que faut-il changer pour détecter en direct une technique aujourd'hui vue
seulement après coup~?
\end{boxexo}

\begin{boxexo}[title={\ding{45}\quad Exercice~21.7. LOGONE \niveaumoyen}]
À partir du scénario, proposez une troisième règle utile pour LiZbiZ.
\end{boxexo}

\begin{boxexo}[title={\ding{45}\quad Exercice~21.8. Livraison \niveaumoyen}]
Pourquoi la liste horodatée des actions du testeur est-elle un bon jeu de
test pour le SOC~?
\end{boxexo}

\begin{boxexo}[title={\ding{45}\quad Exercice~21.9. Couverture complète \niveauexpert}]
Construisez la matrice de couverture des techniques des Chapitres~\ref{chap:reseau}
à~\ref{chap:post}.
\end{boxexo}

\begin{boxexo}[title={\ding{45}\quad Exercice~21.10. Programme \niveauexpert}]
Proposez un programme d'amélioration continue de la détection fondé sur des
exercices violets trimestriels.
\end{boxexo}

\bigskip

\begin{boxinfo}
\textbf{Pour aller plus loin.}
\begin{itemize}
    \item MITRE ATT\&CK~\cite{mitreattack}.
\end{itemize}
\end{boxinfo}

\begin{boxinfo}
\textbf{Transition.} Toutes ces conclusions doivent être restituées~: la
Partie~VI commence par le Chapitre~\ref{chap:rapport}, consacré au rapport.
\end{boxinfo}
